\chapter{Obligations d'État européennes et japonaises}\label{ch:m2:european-and-japanese-government-bonds}

Le mercredi 28 septembre 2022, le rendement de l'obligation d'État britannique à trente ans
a varié de 127 points de base en une seule journée, soit davantage que son
amplitude sur toute une année dans tous les cas sauf quatre des vingt-sept années précédentes. Dans les jours
précédents, il avait progressé de 160 points de base, ayant débuté l'année autour de 1.2~\%.
Les fonds de pension qui avaient emprunté en donnant leurs \omterm{def:m2:european-and-japanese-government-bonds:ldi}{gilts} en garantie pour apparier leurs engagements à long
terme se voyaient réclamer du collatéral plus vite qu'ils ne pouvaient en réunir,
et leurs gérants disaient à la Banque d'Angleterre que, aux rendements en
vigueur, certains de leurs fonds auraient une valeur négative dès le lendemain
matin. Ce jour-là, la Banque annonça des achats temporaires de \omterm{def:m2:european-and-japanese-government-bonds:ldi}{gilts} longs pour
restaurer des conditions de marché ordonnées~; d'ici le 14 octobre, elle en avait acquis pour
19.3~milliards de livres. La zone euro, le Japon et le Royaume-Uni gèrent
leurs marchés d'obligations d'État de façon différente, et chacun a dû, au cours des quinze
dernières années, trancher ce que sa banque centrale doit au marché en cas de crise. Ce
chapitre les compare.

\section{La zone euro~: une monnaie, de nombreux émetteurs}

Vingt et un États empruntent en euros depuis l'adhésion de la Bulgarie le 1\up{er} janvier
2026, et aucun d'eux ne contrôle la banque centrale qui émet cette monnaie. La même monnaie se traite donc à des rendements multiples, et les
écarts traduisent la mesure, par le marché, du crédit de chaque État et du
risque que l'union monétaire elle-même puisse changer.

\begin{definition}[Spread souverain et risque de relibellé]\label{def:m2:european-and-japanese-government-bonds:spread}
Un \emph{spread souverain}\index{spread souverain} est l'écart entre le
rendement d'une obligation d'État et celui d'une obligation d'État de référence de même
devise et de même échéance~; dans la zone euro, il s'agit en général de l'Allemagne. \emph{risque de
relibellé}\index{risque de relibellé} est la part d'un spread qui rémunère la
possibilité qu'une obligation soit remboursée dans une nouvelle monnaie nationale valant moins
que l'euro.
\end{definition}

\begin{definition}[Bund et syndication]\label{def:m2:european-and-japanese-government-bonds:bund}
Un \emph{Bund}\index{Bund} est une obligation d'État fédérale allemande (Bundesanleihe),
émise à 7, 10, 15 ou 30 ans avec un coupon annuel~; le rendement du Bund à dix ans
est le taux long de référence de la zone euro. \emph{syndication}\index{syndication}
est la vente d'une nouvelle obligation par un groupe de banques qui constituent un
carnet d'ordres d'investisseurs et fixent le prix, au lieu d'une adjudication~; certains émetteurs y recourent
pour les nouvelles obligations longues, là où la demande est la plus difficile à prévoir.
\end{definition}

\begin{omfigure}
\begin{tikzpicture}
\begin{axis}[width=11cm,height=5cm,
  ylabel={spread par rapport à l'Allemagne (pb)},
  tick label style={font=\small},label style={font=\small},
  xmin=2007,xmax=2027,ymin=-20,ymax=560,grid=major,grid style={gray!25},
  xtick={2008,2011,2014,2017,2020,2023,2026},
  xticklabel style={/pgf/number format/1000 sep=},
  legend columns=3,legend style={font=\footnotesize,at={(0.5,-0.18)},anchor=north,draw=none,fill=none,/tikz/every even column/.append style={column sep=8pt}},
  table/col sep=comma]
\addplot[omThm,very thick] table[x=t,y=IT]{figdata/markets-2/07-european-and-japanese-government-bonds/spreads.csv};
\addplot[omProp,thick] table[x=t,y=ES]{figdata/markets-2/07-european-and-japanese-government-bonds/spreads.csv};
\addplot[omDef,thick,dashed] table[x=t,y=FR]{figdata/markets-2/07-european-and-japanese-government-bonds/spreads.csv};
\legend{Italie,Espagne,France}
\node[font=\footnotesize,anchor=west] at (axis cs:2014.3,480) {Italie, nov. 2011~: 518 pb};
\end{axis}
\end{tikzpicture}

\omcaption{Spreads à dix ans avec l'Allemagne, moyennes mensuelles, de 2007 à août 2026.
La crise de la dette souveraine a porté celui de l'Italie à 518 points de base en
novembre 2011 et celui de l'Espagne à 555 en juillet 2012~; les spreads se sont resserrés après que la BCE eut annoncé ses opérations monétaires sur titres (OMT) en
septembre 2012. En août 2026, le spread de la France, 82 points de base, était légèrement supérieur à
celui de l'Italie, 81. Données~: OCDE, via FRED.}\label{fig:m2:european-and-japanese-government-bonds:spreads}
\end{omfigure}

\begin{example}[Anatomie d'un élargissement]\label{ex:m2:european-and-japanese-government-bonds:2011}
Entre juin et novembre 2011 (moyennes mensuelles), le spread à dix ans de l'Italie
par rapport à l'Allemagne s'est élargi de 326 points de base. Le rendement italien a augmenté de 224 points de base~; les
autres 102 points proviennent de la \emph{baisse} de même ampleur du rendement allemand, à mesure que les investisseurs se
réfugiaient vers les \omterm{def:m2:european-and-japanese-government-bonds:bund}{Bund}. Un spread mesure le crédit relatif, mais la moitié d'un élargissement
en période de crise peut venir de la remontée des cours de l'indice de référence~: une couverture qui vend cet indice à découvert contre le
spread perd sur cette moitié.
\end{example}

\section{Spreads, filets de sécurité et marché électronique des dealers}

\begin{definition}[Filet de sécurité de la banque centrale]\label{def:m2:european-and-japanese-government-bonds:backstop}
Un \emph{filet de sécurité de la banque centrale}\index{filet de sécurité de la banque centrale} est un engagement, ou
un instrument permanent par lequel une banque centrale achètera les obligations d'un État
en cas de dysfonctionnement du marché, afin que les cours reflètent les fondamentaux plutôt qu'une
ruée. Ses conditions déterminent à quoi il peut servir~: les conditions d'activation,
sa taille, et le fait qu'il soit annoncé ou non comme illimité.
\end{definition}

La BCE en dispose de deux. Les opérations monétaires sur titres (OMT), annoncées en septembre 2012,
n'ont pas de limite ex ante mais sont subordonnées à ce que le pays dispose d'un programme d'ajustement
avec le Mécanisme européen de stabilité~: la promesse est aussi grande que
nécessaire, et son prix est un programme. L'instrument de protection de la transmission
de juillet 2022 est en revanche subordonné à quatre critères de solidité budgétaire
et macroéconomique, et vise des «~dynamiques de marché injustifiées et désordonnées
qui font peser une menace sérieuse sur la transmission de la politique monétaire~»,
des achats concentrés sur les échéances d'un à dix ans. Un spread qui
s'élargit pour des raisons fondamentales reste hors des deux~; un spread qui s'élargit parce que
d'une ruée est, en principe, couvert. Trancher entre les deux revient à la BCE.

Le marché secondaire est organisé à deux niveaux, comme le marché américain. Entre dealers,
les obligations d'État européennes se négocient sur des plateformes électroniques, comme les
marchés inter-dealers de MTS, un par pays~; les clients traitent avec les dealers par
demande de cotation (RFQ) sur des \omterm{def:m2:the-treasury-market:idb}{plateformes dealer-to-client}, comme BondVision, la plateforme propre à MTS
(voir au \cref{ch:m2:how-bonds-trade}).

\begin{dated}{2026-09}{Rendements à dix ans, moyennes d'août 2026}\label{dat:m2:european-and-japanese-government-bonds:yields}
Allemagne 3.18~\%, Espagne 3.63~\%, Italie 3.99~\%, France 4.00~\%~; Japon 2.94~\%~;
Royaume-Uni 4.99~\%. Spreads par rapport à l'Allemagne~: Espagne 45, Italie 81, France 82 points de base.
\end{dated}

\section{Japon~: contrôle de la courbe des taux et sa sortie}

\begin{definition}[Contrôle de la courbe des taux]\label{def:m2:european-and-japanese-government-bonds:ycc}
Le \emph{contrôle de la courbe des taux}\index{contrôle de la courbe des taux} est une politique par laquelle la
banque centrale fixe un objectif pour le rendement d'une obligation d'État à long terme et achète, en
quantité nécessaire, pour le maintenir à ce niveau, en plus de fixer le taux
directeur à court terme.
\end{definition}

La Banque du Japon l'a introduit en septembre 2016, en s'engageant à acheter
des obligations d'État afin que les rendements à dix ans restent «~à peu près au
niveau actuel (autour de zéro pour cent)~». Un objectif de rendement inverse le
lien habituel entre prix et quantité~: la banque centrale ne choisit plus combien
elle achète~: c'est le marché qui le décide, en lui vendant au niveau visé. Lorsque les rendements
montent partout dans le monde, défendre l'objectif revenait à acheter tout ce qui était offert~; le
20 décembre 2022, la Banque a élargi la bande autour de cet objectif, d'un quart de point
à un demi-point de pourcentage et a proposé d'acheter, chaque jour ouvré, des obligations à dix ans à 0.5~\%.
Le 19 mars 2024, elle a déclaré que le \omterm{def:m2:european-and-japanese-government-bonds:ycc}{contrôle de la courbe des taux} et le
taux négatif avaient «~rempli leur rôle~» et elle est revenue au pilotage du seul
taux court. Le rendement à dix ans, proche de zéro pendant l'essentiel de la période de contrôle,
s'est établi en moyenne à 2.94~\% en août 2026 (\cref{fig:m2:european-and-japanese-government-bonds:jgb}).

\begin{omfigure}
\begin{tikzpicture}
\begin{axis}[width=11cm,height=4.8cm,
  ylabel={rendement du JGB à dix ans (\%)},
  tick label style={font=\small},label style={font=\small},
  xmin=2007,xmax=2027,ymin=-0.5,ymax=3.3,grid=major,grid style={gray!25},
  xtick={2008,2011,2014,2017,2020,2023,2026},
  xticklabel style={/pgf/number format/1000 sep=},
  table/col sep=comma]
\fill[omProp!12] (axis cs:2016.72,-0.5) rectangle (axis cs:2024.21,3.3);
\addplot[omDef,very thick] table[x=t,y=JP]{figdata/markets-2/07-european-and-japanese-government-bonds/jgb.csv};
\node[font=\footnotesize,text=omProp,align=center] at (axis cs:2020.45,2.6) {contrôle de la courbe des taux\\sept. 2016 -- mars 2024};
\end{axis}
\end{tikzpicture}

\omcaption{Rendement de l'obligation d'État japonaise à dix ans, moyennes mensuelles, de 2007
à août 2026. Sous le \omterm{def:m2:european-and-japanese-government-bonds:ycc}{contrôle de la courbe des taux} (zone grisée), la Banque le maintenait près de son
objectif~; après la sortie, il a atteint les niveaux les plus élevés de la période. Données~:
OCDE, via FRED~; dates de la politique, d'après les
communiqués de la Banque du Japon.}\label{fig:m2:european-and-japanese-government-bonds:jgb}
\end{omfigure}

\section{Le Royaume-Uni et la crise des gilts de 2022}

\begin{definition}[Gilt et gestion adossée au passif]\label{def:m2:european-and-japanese-government-bonds:ldi}
Un \emph{gilt}\index{gilt} est une obligation d'État britannique~; les gilts classiques versent
des coupons semestriels, tandis que les gilts indexés versent coupons et principal
indexés sur l'inflation. La \emph{gestion adossée au passif}\index{gestion adossée au passif}
(LDI) est une stratégie dans laquelle un fonds de pension à prestations définies
couvre la sensibilité de ses engagements aux taux d'intérêt et à l'inflation grâce
à de longs gilts et à des swaps, souvent avec un effet de levier issu des pensions livrées et des dérivés, de sorte que le fonds
peut couvrir une part de ses engagements supérieure à celle de ses obligations en portefeuille, tout en conservant d'autres actifs,
à rendement plus élevé.
\end{definition}

L'effet de levier est ce qui a rendu la couverture fragile. Un fonds qui détient des \omterm{def:m2:european-and-japanese-government-bonds:ldi}{gilts} valant deux fois
son capital, en empruntant l'autre moitié en pension livrée, perd deux fois plus vite quand les \omterm{def:m2:european-and-japanese-government-bonds:ldi}{gilts}
baissent. Son gérant conserve un coussin de capital face aux pertes et réclame
un apport au fonds de pension lorsque ce coussin s'amenuise~; pour les fonds mutualisés, adossés à de nombreux petits
régimes de retraite, cette demande prend des jours. Fin septembre 2022, les rendements
ont bougé en quelques heures. La seule source rapide de liquidités des fonds était leurs \omterm{def:m2:european-and-japanese-government-bonds:ldi}{gilts}~:
les vendre faisait monter les rendements, ce qui appelait davantage de liquidités et forçait
de nouvelles ventes (\cref{fig:m2:european-and-japanese-government-bonds:loop}). Le 26
septembre, les gestionnaires ont indiqué à la Banque que, à valeur nominale, ils devraient vendre
au moins 50 milliards de livres de longs \omterm{def:m2:european-and-japanese-government-bonds:ldi}{gilts} en peu de temps, alors que les transactions
moyennes atteignaient 12 milliards par jour sur ces échéances.

\begin{omfigure}
\begin{tikzpicture}[font=\small,
  box/.style={draw,thick,rounded corners=2pt,align=center,minimum height=0.9cm,minimum width=3.4cm,inner sep=4pt},
  arr/.style={-{Stealth[length=2.5mm]},thick}]
\node[box,draw=omThm,fill=omThm!8] (y) at (0,1.8) {hausse des rendements des longs gilts};
\node[box,draw=omDef,fill=omDef!8] (c) at (5.6,1.8) {les coussins LDI s'amenuisent};
\node[box,draw=omProp,fill=omProp!8] (m) at (5.6,-0.6) {collatéral et appels de marge\\sur pensions livrées et swaps};
\node[box,draw=gray,fill=gray!8] (s) at (0,-0.6) {ventes forcées de longs gilts\\sur un marché étroit};
\draw[arr] (y) -- node[above=11pt] {les prix baissent} (c);
\draw[arr] (c) -- node[right=4pt] {l'effet de levier augmente} (m);
\draw[arr] (m) -- node[below=17pt] {liquidités requises immédiatement} (s);
\draw[arr] (s) -- node[left=4pt] {offre accrue} (y);
\node[font=\footnotesize,text=omThm,align=center] at (2.8,0.6) {la boucle que la Banque\\a brisée le 28 septembre};
\end{tikzpicture}

\omcaption{La spirale du collatéral de septembre 2022. Les stratégies LDI des fonds de pension
s'étaient financées en empruntant contre de longs \omterm{def:m2:european-and-japanese-government-bonds:ldi}{gilts}~; la hausse des rendements a érodé leurs coussins, les appels
de collatéral pouvaient être satisfaits au plus vite par la vente de \omterm{def:m2:european-and-japanese-government-bonds:ldi}{gilts}, et ces ventes ont fait
monter encore les rendements. Un acheteur capable d'absorber les ventes sans chercher
de rendement, à savoir la banque centrale, a mis fin à la boucle.}\label{fig:m2:european-and-japanese-government-bonds:loop}
\end{omfigure}

Les achats de la Banque ont été ciblés et temporaires~: de longs \omterm{def:m2:european-and-japanese-government-bonds:ldi}{gilts} uniquement, du 28
septembre au 14 octobre, 12.1 milliards de \omterm{def:m2:european-and-japanese-government-bonds:ldi}{gilts} classiques et 7.2 milliards de \omterm{def:m2:european-and-japanese-government-bonds:ldi}{gilts}
indexés, revendus au marché entre novembre 2022 et janvier
2023. Le Comité de politique financière a alors fixé une norme~: les fonds LDI devraient
être en mesure de résister à une hausse d'environ 250 points de base des rendements des \omterm{def:m2:european-and-japanese-government-bonds:ldi}{gilts} sans
ventes forcées, contre les 100 points de base que les tests de résistance précédents
avaient retenus.

\begin{proposition}[Jusqu'où peut chuter un fonds obligataire à effet de levier]\label{prop:m2:european-and-japanese-government-bonds:cushion}
Un fonds détient des obligations valant $G$, de \omterm{def:m2:government-bonds:duration}{duration modifiée} $D$ et de convexité $\mathcal C$,
financées par une pension livrée $B = \ell G$, de sorte que son coussin vaut $(1-\ell)G$. Son coussin est
épuisé par la hausse parallèle des rendements $\Delta y^*$ à partir de laquelle la valeur des obligations tombe
à $B$~; au premier ordre $\Delta y^* \approx (1-\ell)/D$. La perte du second ordre
$D\,\Delta y - \frac12\,\mathcal C\,(\Delta y)^2$ ne dépasse jamais $D^2/(2\mathcal C)$, de sorte que
lorsque $1 - \ell > D^2/(2\mathcal C)$, le trinôme n'a pas de racine et seule une réévaluation complète
donne $\Delta y^*$. Après une hausse laissant les obligations valant $G'$, le
fonds doit vendre $2B - G'$ d'entre elles pour rétablir une part de pension livrée égale à un demi, sans
nouveau capital.
\end{proposition}
\begin{proof}
Le coussin vaut $G' - B$. Au premier ordre $G'/G = 1 - D\Delta y$~; en posant $G' = B$
donne la première affirmation. Le trinôme $D x - \frac12 \mathcal C x^2$ est maximal en
$x = D/\mathcal C$, avec pour valeur $D^2/(2\mathcal C)$. Après la vente de $S$ et le remboursement
de $S$ de pension livrée, la part vaut $(B - S)/(G' - S) = \frac12$, donc $S = 2B - G'$.
\end{proof}

Pour l'obligation d'État à long terme de la \cref{fig:m2:european-and-japanese-government-bonds:ldi},
$D = 21.8$ et $\mathcal C = 589$~: la perte maximale du trinôme vaut 40.5~\%, soit moins
que le coussin de 50~\% du fonds, si bien que la formule du second ordre laisserait entendre que le coussin
ne peut jamais être épuisé. Il l'est à 351 points de base. Les développements de Taylor
valent pour les petits mouvements~; les tests de résistance portent précisément sur les grands mouvements.

\begin{omfigure}
\begin{tikzpicture}
\begin{axis}[width=11cm,height=4.6cm,
  xlabel={hausse des rendements des gilts (pb)},ylabel={coussin restant (\%)},
  tick label style={font=\small},label style={font=\small},
  xmin=0,xmax=400,ymin=-25,ymax=105,grid=major,grid style={gray!25},
  table/col sep=comma]
\addplot[omDef,very thick] table[x=shock,y=cushion]{figdata/markets-2/07-european-and-japanese-government-bonds/ldi.csv};
\draw[omThm,thick,dashed] (axis cs:160,-25) -- (axis cs:160,105);
\draw[omProp,thick,dashed] (axis cs:250,-25) -- (axis cs:250,105);
\node[font=\footnotesize,anchor=south east,text=omThm,align=right] at (axis cs:156,-22) {Septembre 2022~:\\+160};
\node[font=\footnotesize,anchor=south west,text=omProp] at (axis cs:253,55) {standard~: 250};
\end{axis}
\end{tikzpicture}

\omcaption{Le coussin d'un fonds LDI illustratif, avec un effet de levier de deux en pension livrée
et détenant une obligation d'État à trente ans, lorsque les rendements des \omterm{def:m2:european-and-japanese-government-bonds:ldi}{gilts} passent de 3.5~\%. Il tombe à 43~\%
de sa valeur initiale après 160 points de base, l'ampleur du mouvement de septembre 2022,
et à zéro après 351. Données~: le problème du
week-end du chapitre.}\label{fig:m2:european-and-japanese-government-bonds:ldi}
\end{omfigure}

\begin{dated}{2026-09}{La détention de gilts par la Banque d'Angleterre}\label{dat:m2:european-and-japanese-government-bonds:uk}
Le 17 septembre 2026, la Banque a suspendu ses ventes aux enchères de \omterm{def:m2:european-and-japanese-government-bonds:ldi}{gilts} issus de l'Asset
Purchase Facility, le temps d'examiner la suite à donner~: les \omterm{def:m2:european-and-japanese-government-bonds:ldi}{gilts} arrivant à échéance avant 2035
(222~milliards de livres) seront conservés jusqu'à l'échéance, les plus longs (120~milliards de livres) conservés
jusqu'à l'échéance pour garantir l'émission de billets, et ceux arrivant à échéance entre 2035 et 2049
(146~milliards de livres à l'achat) pourraient être vendus au gouvernement à raison de 20~milliards de livres par
an, sous réserve d'un examen avant avril 2027.
\end{dated}

\section{Tutoriel~: un moniteur de spreads souverains}\label{tut:m2:european-and-japanese-government-bonds:monitor}

\begin{tutorial}
\textbf{Objectif.} Construire les spreads de la zone euro à partir des rendements mensuels, décomposer l'élargissement de 2011
et signaler les mouvements inhabituels. \textbf{État final~:}
\cref{fig:m2:european-and-japanese-government-bonds:spreads} et les résultats de
l'\cref{ex:m2:european-and-japanese-government-bonds:2011}.
\begin{tutsteps}
\item \textbf{Spreads et leur décomposition.}
\omcode{code/firm/sovspread/firm_sovspread.py}{10}{31}{Spreads par rapport à un indice de référence, et variation d'un spread décomposée entre les mouvements des deux rendements.}\label{lst:m2:european-and-japanese-government-bonds:spreads}
\item \textbf{Une alerte z-score} sur une fenêtre glissante.
\omcode{code/firm/sovspread/firm_sovspread.py}{34}{42}{Z-score glissant d'un spread.}\label{lst:m2:european-and-japanese-government-bonds:z}
\item \textbf{Exécuter} \texttt{sovereign\_demo.widening\_2011()}~: 326, 224 et 102
points de base~; \texttt{fig\_sovereign.py} produit les trois graphiques.
\item \textbf{Le fonds LDI.}
\omcode{code/markets-2/07-european-and-japanese-government-bonds/python/sovereign_demo.py}{51}{59}{Gilts, pension livrée et coussin d'un fonds LDI après une hausse des rendements.}\label{lst:m2:european-and-japanese-government-bonds:fund}
\end{tutsteps}
\textbf{Pour aller plus loin.} Appliquer l'alerte z-score au spread italien avec une
fenêtre de 24~mois et lister les mois signalés~; puis donner au fonds LDI un effet de levier de trois
et trouver la hausse des rendements qui l'épuise.
\end{tutorial}

\section{Construction~: le moniteur de spreads souverains}\label{bld:m2:european-and-japanese-government-bonds:monitor}

\begin{build}
\bfield{Rôle} La firme miniature a un desk de taux qui négocie des spreads de la zone euro et
a besoin d'une vue en temps réel de ces spreads, de ce qui a fait bouger chacun d'eux, et de savoir quels mouvements sont
inhabituels.
\bfield{Interface} \texttt{spreads(yields, benchmark)}~;
\texttt{decompose(issuer, benchmark, start, end)} renvoyant
\texttt{Decomposition(change\_bp, issuer\_bp, benchmark\_bp)}~;
\texttt{zscore(series, window)}~; \texttt{alerts(series, window, threshold)}.
\bfield{Règles} Rendements en pourcentage, spreads en points de base~; spreads uniquement aux
dates communes aux deux émetteurs~; le z-score d'une valeur utilise la fenêtre
\emph{précédant} celle-ci.
\bfield{Tests d'acceptation} \texttt{code/firm/sovspread/tests/}~: spreads aux dates communes~; la décomposition
s'additionne et donne à la baisse de l'indice de référence un
signe d'élargissement~; un saut déclenche l'alerte et le bruit ne la déclenche pas.
\bfield{Pour aller plus loin} Données quotidiennes issues des courbes de référence des émetteurs eux-mêmes~; spreads
à échéances appariées avec la bibliothèque d'obligations
(\cref{bld:m2:government-bonds:bond})~; un spread en termes d'asset swap
(\cref{ch:m2:corporate-bonds}).
\end{build}

\begin{omsources}
\item Bank of England, lettre de Sir Jon Cunliffe à la Treasury Committee sur
les LDI, 5 octobre 2022~; \emph{Quarterly Bulletin} 2023, ``Financial stability buy/sell
tools: a gilt market case study''~; document de travail sur la résilience minimale des LDI, 2023~;
avis de marché de l'APF, 17 septembre 2026.
\item European Central Bank, communiqués de presse sur les outright monetary transactions
(6 septembre 2012) et sur le transmission protection instrument (21 juillet 2022).
\item Bank of Japan, ``New Framework for Strengthening Monetary Easing'', 21
septembre 2016~; ``Changes in the Monetary Policy Framework'', 19 mars 2024.
\item OECD, long-term interest rates (via FRED)~; MTS, site de l'entreprise.
\end{omsources}

\section{Exercices}

\begin{exercise}[$\star$]\label{exo:m2:european-and-japanese-government-bonds:1}
Avec les rendements de la \cref{dat:m2:european-and-japanese-government-bonds:yields},
calculer les spreads de l'Espagne, de l'Italie et de la France par rapport à l'Allemagne, en points de base.
\end{exercise}

\begin{exercise}[$\star$]\label{exo:m2:european-and-japanese-government-bonds:2}
Le rendement italien monte de 30 points de base et le rendement allemand baisse de 10 en une journée. Calculer
la variation du spread et sa décomposition.
\end{exercise}

\begin{exercise}[$\star$]\label{exo:m2:european-and-japanese-government-bonds:3}
Un fonds LDI détient 1~milliard de livres de \omterm{def:m2:european-and-japanese-government-bonds:ldi}{gilts}, financés à moitié en pension livrée. Calculer son
effet de levier. Calculer son coussin après une baisse de 20~\% des cours des \omterm{def:m2:european-and-japanese-government-bonds:ldi}{gilts}, ainsi que son nouvel
effet de levier.
\end{exercise}

\begin{exercise}[$\star\star$]\label{exo:m2:european-and-japanese-government-bonds:4}
En utilisant la \cref{prop:m2:european-and-japanese-government-bonds:cushion}, estimer la
hausse de rendement qui épuise le coussin d'un fonds de duration 21.8 et de convexité 589,
dont la moitié des \omterm{def:m2:european-and-japanese-government-bonds:ldi}{gilts} est en pension livrée, au premier ordre. Que dit la formule au second ordre
et pourquoi~?
\end{exercise}

\begin{exercise}[$\star\star$]\label{exo:m2:european-and-japanese-government-bonds:5}
Expliquer pourquoi la Banque d'Angleterre n'a acheté que des \omterm{def:m2:european-and-japanese-government-bonds:ldi}{gilts} longs en 2022, et seulement
pendant une quinzaine de jours. Comparer avec l'\omterm{def:m2:central-banks-and-the-short-rate:qe}{assouplissement quantitatif}.
\end{exercise}

\begin{exercise}[$\star\star$]\label{exo:m2:european-and-japanese-government-bonds:6}
Dans le cadre du \omterm{def:m2:european-and-japanese-government-bonds:ycc}{contrôle de la courbe des taux}, le marché vend des obligations à dix ans
au rendement cible. Expliquer qui décide du nombre d'obligations que la banque centrale achète
et ce qu'il advient de son bilan si le marché anticipe l'abandon de la cible.
\end{exercise}

\begin{exercise}[$\star\star\star$]\label{exo:m2:european-and-japanese-government-bonds:7}
\emph{Programmation.} Avec \texttt{sovereign\_demo}, calculer les ventes que le fonds LDI
doit réaliser pour retrouver un levier de moitié après des hausses de 100 et de 160 points
de base, ainsi que le coussin (en part de la détention) nécessaire pour survivre à 250.
\end{exercise}

\begin{exercise}[$\star\star\star$]\label{exo:m2:european-and-japanese-government-bonds:8}
\emph{Trouver l'erreur.} «~Les fonds LDI ont perdu de l'argent en 2022 parce que la hausse des rendements
a appauvri les régimes de retraite.~» Corriger cette affirmation en s'appuyant sur
l'exposé propre à la Banque d'Angleterre sur les bilans des régimes.
\end{exercise}

\section{Problème~: la spirale du collatéral}

\begin{problem}[{Problème du week-end --- jusqu'où peut chuter un fonds LDI}]\label{pb:m2:european-and-japanese-government-bonds:1}
Le 22 septembre 2022, un fonds LDI détient 1~milliard de livres (valeur de marché)
d'un \omterm{def:m2:european-and-japanese-government-bonds:ldi}{gilt} à trente ans, à coupon de 1.5~\%, échéant le 22 juillet 2052, à un rendement de
3.50~\%~; la moitié est financée en pension livrée. Le fonds ne peut pas lever de nouveaux fonds propres en moins
d'une semaine. Le \omterm{def:m2:european-and-japanese-government-bonds:ldi}{gilt} et le fonds sont donnés à titre d'illustration.

\medskip\noindent\textbf{Partie I --- Le fonds.}
\begin{enumerate}
\item Donner le prix plein coupon du \omterm{def:m2:european-and-japanese-government-bonds:ldi}{gilt}, la pension livrée du fonds et son coussin.
\item Donner la \omterm{def:m2:government-bonds:duration}{duration modifiée} du \omterm{def:m2:european-and-japanese-government-bonds:ldi}{gilt} et le DV01 du fonds par point de base.
\item Calculer l'effet de levier du fonds, ainsi que son DV01 par livre de fonds propres.
\item Calculer, au premier ordre, la hausse en points de base qui coûte 1~\% du coussin.
\item Expliquer pourquoi un régime de retraite choisirait un tel fonds plutôt que de détenir
directement des \omterm{def:m2:european-and-japanese-government-bonds:ldi}{gilts}.
\end{enumerate}

\medskip\noindent\textbf{Partie II --- Le choc.}
\begin{enumerate}[resume]
\item Les rendements montent de 100 points de base. Donner la valeur des \omterm{def:m2:european-and-japanese-government-bonds:ldi}{gilts}, le coussin et
l'effet de levier.
\item Même calcul pour 160 points de base, l'ampleur du mouvement de septembre.
\item Calculer les ventes que doit réaliser le fonds après 100 points de base pour retrouver
un levier de moitié sans nouveaux fonds propres~? Et après 160~?
\item Comparer les ventes globales déclarées par les gérants avec le volume quotidien
du marché, puis expliquer ce qu'il advient du prix des actifs vendus.
\item Expliquer pourquoi les appels de marge ont porté sur les swaps autant que sur la pension livrée.
\end{enumerate}

\medskip\noindent\textbf{Partie III --- La limite.}
\begin{enumerate}[resume]
\item Donner la hausse qui épuise le coussin.
\item Expliquer ce qu'il advient du prêteur de la pension livrée lorsque le coussin est épuisé.
\item Donner le coussin dont le fonds a besoin pour survivre à 250 points de base sans
vendre.
\item Les tests de résistance antérieurs retenaient 100 points de base. Calculer le coussin que cela impliquait.
\item Expliquer pourquoi un coussin exprimé en pourcentage ne suffit pas à établir la résilience sans
la duration de ce qui est détenu.
\end{enumerate}

\medskip\noindent\textbf{Partie IV --- Jugement.}
\begin{enumerate}[resume]
\item Expliquer pourquoi une banque centrale confrontée à un problème d'inflation a acheté des obligations longues.
\item Expliquer pourquoi les achats ont été temporaires et modestes au regard du marché.
\item Expliquer pourquoi le régime de retraite lui-même, et non le fonds, n'était pas insolvable.
\item Énoncer le \emph{résultat nommé}~: la hausse de rendement qui épuise le coussin,
et le coussin après le mouvement de septembre.
\item En une phrase, dire ce qui a transformé une hausse des rendements en crise.
\end{enumerate}
\end{problem}

\section{Questions d'entretien}

\begin{interviewq}[$\star$ \iqroles{trader, chercheur}]\label{iq:m2:european-and-japanese-government-bonds:1}
Pourquoi les obligations italiennes et allemandes à dix ans, libellées dans la même devise,
se négocient-elles à des rendements différents~? Qu'y a-t-il dans le spread~?
\end{interviewq}

\begin{interviewq}[$\star$ \iqroles{trader, banque}]\label{iq:m2:european-and-japanese-government-bonds:2}
Que s'est-il passé sur le marché des \omterm{def:m2:european-and-japanese-government-bonds:ldi}{gilts} en septembre 2022~?
\end{interviewq}

\begin{interviewq}[$\star\star$ \iqroles{trader, chercheur}]\label{iq:m2:european-and-japanese-government-bonds:3}
Comment fonctionne le \omterm{def:m2:european-and-japanese-government-bonds:ycc}{contrôle de la courbe des taux}, et quels sont les risques d'en sortir~?
\end{interviewq}

\begin{interviewq}[$\star\star$ \iqroles{chercheur, trader}]\label{iq:m2:european-and-japanese-government-bonds:4}
Un spread s'est élargi de 100 points de base. Comment distinguez-vous la part revenant à l'émetteur
et celle revenant à l'indice de référence, et pourquoi cette distinction importe-t-elle pour une couverture~?
\end{interviewq}

\begin{interviewq}[$\star\star$ \iqroles{banque, chercheur}]\label{iq:m2:european-and-japanese-government-bonds:5}
Quelle est la différence entre les programmes OMT et TPI de la BCE~?
\end{interviewq}

\begin{interviewq}[$\star\star\star$ \iqroles{chercheur, banque}]\label{iq:m2:european-and-japanese-government-bonds:6}
On vous demande de soumettre à un test de résistance un fonds obligataire à effet de levier. Comment choisissez-vous
l'ampleur du choc, et que doit encore saisir le test~?
\end{interviewq}
